\chapter{FX Connectivity}\label{ch:nw:fx-connectivity}

On EBS each currency pair now matches in one place, New York or London, and a trader in Tokyo reaches it through a gateway
that forwards the order across an ocean or a continent. A liquidity provider in London that streams prices to Tokyo quotes on
information that is two Eurasian crossings old by the time a Tokyo taker sees it: 141 milliseconds at a specialist's \omterm{def:nw:the-north-american-map:factor}{route
factor}. Yet the provider's price check needs no hold to catch a Tokyo taker who saw the news first, as long as the provider
learns of the news as fast as the taker's request travels. When it learns more slowly, the hold that would close the gap is a
delay the market's code of conduct does not allow.

Book~2 described the FX market's structure: primary venues, ECNs, single-dealer platforms, prime brokerage, streaming quotes,
last look and hold times. Book~10 introduced the liquidity aggregator and Book~13 the FIX session. This chapter puts them on the
map of chapters~10 to~12: where FX matching happens, where credit is checked, how a liquidity provider's streams reach the
sites, and what aggregators see. The venues' sites are dated rows; the staleness model is labelled as such.

\section{The matching sites}

\begin{definition}[Multi-site matching]\label{def:nw:fx-connectivity:multisite}
\emph{Multi-site matching}\index{multi-site matching} is a venue design in which matching engines run in several data centres,
each instrument is matched at exactly one of them, and gateways at the other sites forward orders to the instrument's engine
rather than match them locally.
\end{definition}

\begin{dated}{2026-09}{Where FX venues match and stream, from their connectivity documents}\label{dat:nw:fx-connectivity:sites}
\begin{itemize}
\item EBS Market on CME Globex: each product trades in a single location. New York FX spot and metals in Equinix Secaucus NY5
  (backup in Slough LD4.2); London FX spot and NDFs in LD4.2 (backup in NY5). Tokyo TY3 has only a convenience gateway, routing
  orders to New York or London, and credit-screened conflated market data; an order sent to a segment gateway for an instrument
  not matched there is rejected. Shared services, including global credit limits, in Chicago DC3.
\item LSEG FX (guide of June 2026): Spot and Forwards Matching in London LD4 (recovery in NY6); PriceStream streams and
  aggregation in New Jersey NY6, London LD4 and Tokyo TY3; NDF Matching in Singapore SG3 (recovery in LD4); \omterm{def:nw:colocation-products-and-how-they-are-sold:colo}{colocation}
  \omterm{def:nw:colocation-products-and-how-they-are-sold:mmr}{cross-connects} at LD4, TY3, NY4/NY5 and SG1, with equal cable lengths within the main buildings.
\end{itemize}
\end{dated}

FX has no single exchange per pair: the same pair is quoted on several venues, each with its own sites, and every venue that
matches a pair matches it in one place. The map is therefore short: Slough, Secaucus, Tokyo and Singapore
(\cref{tab:nw:fx-connectivity:floors}). London LD4 is where the London engines are and where the London liquidity providers
sit; the Secaucus campus hosts EBS's New York engine and LSEG's New Jersey streams; TY3 is an access and streaming site with no
EBS matching; SG1 and SG3 serve LSEG's NDFs.

\begin{table}[ht]
\centering\small
\begin{tabular}{@{}lrrr@{}}
\toprule
From London LD4 to & Fibre floor, one way (ms) & At 1.5 (ms) & Round trip at 1.5 (ms) \\
\midrule
Secaucus NY5 and NY6 & 27.07 & 40.6 & 81.2 \\
Tokyo TY3 & 46.85 & 70.3 & 140.5 \\
Singapore SG1 & 53.10 & 79.7 & 159.3 \\
\bottomrule
\end{tabular}
\caption{One-way fibre floors from London LD4 to the other FX sites (chapter~10's geodesic and \omterm{def:nw:the-north-american-map:floor}{group index}, site coordinates in
\texttt{data/networks/sites.csv}), and at chapter~19's specialist \omterm{def:nw:the-north-american-map:factor}{route factor} of 1.5. Data: \texttt{firm\_fxlinks.one\_way\_ms}.}\label{tab:nw:fx-connectivity:floors}
\end{table}

A Tokyo trader's order to the EBS engines crosses the Tokyo gateway and the network: at a \omterm{def:nw:the-north-american-map:factor}{route factor} of 1.5 and
\qty{40}{\micro\second} of gateway, credit screen and matching stages (assumptions), \qty{79.5}{\milli\second} to New York and
\qty{70.3}{\milli\second} to London, one way.

\section{Credit checks in the order path}

\begin{definition}[Bilateral credit screen]\label{def:nw:fx-connectivity:credit}
A \emph{bilateral credit screen}\index{bilateral credit screen} is a venue's filter that lets two participants trade, and shows
each only the other's prices, when both have granted each other credit; it applies in the market data a participant receives and
in the matching of its orders.
\end{definition}

FX venues are mostly bilateral underneath: a trade is a contract between two named firms, settled between them or their prime
brokers, so each must have granted the other credit. EBS shows each firm a credit-screened book, the price levels it could
actually trade, beside a ``market best'' that ignores credit, and warns that the screened quantity may exceed what is tradeable.
The credit check is therefore in two places on the path: in the data (what the firm is shown) and in matching (what it can
trade), and a firm's apparent best price may not be its best tradeable one.

For a liquidity provider, credit is also the first half of last look. The FX Global Code's Principle 17 names two checks, validity,
including ``sufficient available credit'', and price; and the Global Foreign Exchange Committee's 2021 report adds that ``any delay
that is additional to what is required to complete the price and validity check is considered contrary to Principle 17''
(\cref{dat:nw:fx-connectivity:lastlook}).

\begin{dated}{2026-09}{Last look, from the FX Global Code and its committee}\label{dat:nw:fx-connectivity:lastlook}
Principle 17: last look is a risk control to check validity (operational details and sufficient available credit) and price (whether
the requested price ``remains consistent with the current price that would be available to the Client''). GFXC report on last look
(August 2021): an additional delay beyond the checks is contrary to Principle 17; liquidity providers should disclose their last
look window's maximum and minimum length.
\end{dated}

\section{Integrating as a liquidity provider}

A liquidity provider in London learns of a price move at its origin, prices it and streams the new quote to every venue site. A taker
at a venue site learns of the same move by its own path. The quote's staleness at the site is the provider's lag less the
taker's; a taker who hits the stale quote sends a request back to the provider, which checks the price when the request arrives,
or after a hold.

\omcode{code/firm/fxlinks/firm_fxlinks.py}{72}{82}{Staleness at a venue site, and the hold before the price check that would let the provider know what the taker knew.}\label{lst:nw:fx-connectivity:stale}

\begin{proposition}[When no hold is needed]\label{prop:nw:fx-connectivity:hold}
If the provider's information path from the origin is no slower than the taker's path from the origin followed by the request's
path to the provider, the provider's price check at arrival already reflects every move the taker could have seen; the hold that
closes the gap is the provider's lag less that sum, and zero otherwise.
\end{proposition}

\begin{proof}
A move at time $\tau$ reaches the taker at $\tau + d_T$ and, through the request, the provider no earlier than $\tau + d_T + d_R$;
the provider knows of it at $\tau + d_I + p$. The check at arrival plus a hold $h$ reflects the move if and only if
$d_I + p \le d_T + d_R + h$.
\end{proof}

With every path at a \omterm{def:nw:the-north-american-map:factor}{route factor} of 1.5, the triangle inequality makes the hold the provider's processing alone, whatever the
staleness. For a move in Tokyo, the London provider's quote reaches the Tokyo site \qty{140.6}{\milli\second} after the Tokyo taker
knew of it; the taker's request then takes \qty{70.3}{\milli\second} back to London, where the provider learned of the move
\qty{70.3}{\milli\second} after it happened. The price check at arrival catches it (with the model's jitter of
\qty{0.5}{\milli\second} a leg, 70\% at once and 99.4\% after a two-millisecond hold). If instead the provider learns of Tokyo moves
over a backbone at 2.46 times the floor, its quote is \qty{185.6}{\milli\second} stale at the Tokyo site and the hold that would
close the gap is \qty{45.1}{\milli\second}: the difference between the two \omterm{def:nw:the-north-american-map:factor}{route factors} times the floor.

\omcode{code/networks/24-fx-connectivity/python/nw_fx.py}{33}{38}{The weekend problem's numbers: a Tokyo move, the Tokyo site, a London provider.}\label{lst:nw:fx-connectivity:tokyo}

\begin{omfigure}
\begin{tikzpicture}
\begin{axis}[ybar,width=11.5cm,height=5.8cm,bar width=9pt,ymin=0,ymax=160,xtick=data,
  xticklabels from table={figdata/networks/24-fx-connectivity/staleness.csv}{site},
  ylabel={staleness at the site (ms)},xlabel={venue site},
  nodes near coords={\pgfmathprintnumber[fixed,precision=0]{\pgfplotspointmeta}},
  every node near coord/.append style={font=\scriptsize},tick label style={font=\footnotesize},label style={font=\small},
  grid=major,grid style={gray!25},legend pos=north west,legend style={font=\footnotesize},table/col sep=comma]
\addplot[fill=omThm!45,draw=omThm] table[x=x,y=ny]{figdata/networks/24-fx-connectivity/staleness.csv};
\addplot[fill=omDef!55,draw=omDef] table[x=x,y=tokyo]{figdata/networks/24-fx-connectivity/staleness.csv};
\legend{move in New York,move in Tokyo}
\end{axis}
\end{tikzpicture}

\omcaption{Model: how stale a London liquidity provider's quote is at each FX site (London LD4, Secaucus NY5 and NY6, Tokyo TY3, Singapore SG1), after a move in New York or in Tokyo, every path at
1.5 times the fibre floor and \qty{0.1}{\milli\second} of pricing. A move in London leaves every site \qty{0.1}{\milli\second}
stale. Data: \texttt{fig\_fx.py}, \texttt{nw\_fx.staleness\_table()}.}\label{fig:nw:fx-connectivity:staleness}
\end{omfigure}

\begin{omfigure}
\begin{tikzpicture}
\begin{axis}[width=11cm,height=5.6cm,xmin=0,xmax=60,ymin=0,ymax=1.05,xlabel={hold before the price check (ms)},
  ylabel={informed requests caught},tick label style={font=\footnotesize},label style={font=\small},grid=major,
  grid style={gray!25},legend pos=south east,legend style={font=\footnotesize},table/col sep=comma]
\addplot[omDef,thick] table[x=hold,y=equal]{figdata/networks/24-fx-connectivity/caught.csv};
\addplot[omThm,thick] table[x=hold,y=backbone]{figdata/networks/24-fx-connectivity/caught.csv};
\legend{information path at 1.5,information path at 2.46}
\end{axis}
\end{tikzpicture}

\omcaption{Model: the share of trade requests from a Tokyo taker who saw a Tokyo move first that the London provider's price check
catches, against the hold before the check, with exponential jitter of \qty{0.5}{\milli\second} on each leg. With an information path
as fast as the request's, a two-millisecond hold catches 99.4\%; with a backbone, 47 milliseconds are needed. Data:
\texttt{fig\_fx.py}, \texttt{nw\_fx.caught\_curves()}.}\label{fig:nw:fx-connectivity:caught}
\end{omfigure}

The Code settles which fix is allowed. A hold beyond what the checks require is contrary to Principle 17, so the provider cannot buy
back a slow information path with a longer last look window; it must make the path faster, or price Tokyo moves in Tokyo, which is
what streaming sites such as LSEG's PriceStream in TY3 allow. Staleness itself is not the provider's to remove: a quote made in
London is a round trip old in Tokyo, and the only remedy for a taker who wants fresh London prices in Tokyo is a provider that quotes
from Tokyo.

\section{Aggregators}

An aggregator combines several providers' streams and shows the best. After a move, the providers whose quotes arrive last are the
ones whose prices have not moved, and their stale prices are the best on the side the market moved away from: the aggregator selects
the stalest quote exactly when it is stale. At an aggregator in Tokyo, after a Tokyo move, a Tokyo provider's quote is fresh, a London
provider's stays stale for \qty{140.5}{\milli\second} and a New Jersey provider's for \qty{158.9}{\milli\second}
(\texttt{firm\_fxlinks.stale\_window\_ms}). This is the provider's side of chapter~19's races: a provider
is picked off in proportion to its distance from the moves of the pairs it quotes, and an aggregator makes the picking-off
systematic. The provider's defences are the ones above: a fast information path, pricing near the origin, and a price check without
added delay.

\begin{method}[Connecting as an FX liquidity provider]\label{met:nw:fx-connectivity:lp}
\begin{enumerate}
\item List the venues and their sites for each pair: where it matches, where it streams, where the gateways are.
\item For each pair, find where its moves originate, and compare the provider's information path with the takers' paths and the
  request path back (\cref{prop:nw:fx-connectivity:hold}).
\item Where the provider's information path is slower, shorten it or price from a site near the origin; never lengthen the last look
  window to compensate.
\item Keep credit consistent across venues and in the validity check, and measure rejects by reason.
\item Measure staleness at each site from your own timestamps, and the fills that follow moves you had not yet seen.
\end{enumerate}
\end{method}

\begin{omfigure}
\begin{tikzpicture}[font=\small,
  site/.style={draw=omDef,fill=omDef!10,rounded corners,align=center,minimum height=10mm,inner sep=3pt},
  lp/.style={draw=omThm,fill=omThm!12,rounded corners,align=center,minimum height=10mm,inner sep=3pt}]
\node[lp] (ld) at (0,0) {London LD4\\provider; engines};
\node[site] (ny) at (-4.2,-2.2) {Secaucus NY5 and NY6\\engine; streams};
\node[site] (ty) at (4.2,-2.2) {Tokyo TY3\\gateway; streams};
\node[site] (sg) at (4.2,1.6) {Singapore SG1 and SG3\\NDF engine};
\draw[->,omThm,thick] (ld) -- node[above,sloped,font=\footnotesize] {40.6 ms} (ny);
\draw[->,omThm,thick] (ld) -- node[above,sloped,font=\footnotesize] {70.3 ms} (ty);
\draw[->,omThm,thick] (ld) -- node[above,sloped,font=\footnotesize] {79.7 ms} (sg);
\draw[->,omDef,dashed] (ty) -- node[below,font=\footnotesize] {orders to New York or London} (ny);
\end{tikzpicture}

\omcaption{A London provider's streams to the FX sites, one way at 1.5 times the fibre floor, and the Tokyo gateway that forwards
orders to the engines. Schematic; distances as in \cref{tab:nw:fx-connectivity:floors}.}\label{fig:nw:fx-connectivity:fanout}
\end{omfigure}

\section{Tutorial: staleness at every site}\label{tut:nw:fx-connectivity:tut}

\begin{tutorial}
\textbf{Goal.} Compute how stale a London provider's quotes are at the FX sites, and what its price check needs to catch the takers
who saw the move first. \textbf{End state:} \cref{fig:nw:fx-connectivity:staleness,fig:nw:fx-connectivity:caught}.
\begin{tutsteps}
\item \textbf{Sites.} \texttt{firm\_fxlinks.load\_venues} and \texttt{sites} read the dated venue rows and the coordinates.
\item \textbf{Staleness.} \texttt{staleness\_ms} and \texttt{hold\_ms} (\cref{lst:nw:fx-connectivity:stale}) for every origin and
  site, with \texttt{Paths} for the \omterm{def:nw:the-north-american-map:factor}{route factors}.
\item \textbf{Checks.} \texttt{caught} adds jitter and counts the informed requests caught after each hold; \texttt{nw\_fx.tokyo}
  (\cref{lst:nw:fx-connectivity:tokyo}) gives the problem's numbers.
\end{tutsteps}
\textbf{What to change next.} Give each taker a microwave path where one exists (chapter~14) and see which sites' holds jump; add a
second provider in Tokyo and measure what an aggregator there selects.
\end{tutorial}

\section{Build: the FX connectivity model}\label{bld:nw:fx-connectivity:fxlinks}

\begin{build}
\bfield{Purpose} Where FX venues match and stream, and what a provider's streams and price checks see at each site: the FX rows of
chapter~29's plan.
\bfield{Interface} \texttt{firm\_fxlinks}: \texttt{Venue}, \texttt{load\_venues}, \texttt{sites}, \texttt{one\_way\_ms},
\texttt{Paths}, \texttt{staleness\_ms}, \texttt{hold\_ms}, \texttt{caught}, \texttt{stale\_window\_ms}, \texttt{order\_path\_ms}.
\bfield{Rules} Venue sites are dated rows with sources; paths are fibre floors times stated \omterm{def:nw:the-north-american-map:factor}{route factors}; jitter and processing are
stated assumptions; results are a model.
\bfield{Acceptance tests} \texttt{code/firm/fxlinks/tests/}: the venue table, no hold with equal paths for every origin and site, the
hold equal to the route-factor gap, and the stale windows and order path.
\bfield{Stretch} Microwave paths for takers; several providers at several sites; credit limits per counterparty in the order path.
\end{build}

\begin{omsources}
\item CME Group Client Systems Wiki: EBS Market on CME Globex Market Functionality; iLink Binary Order Entry for EBS; credit-screened
  market data for EBS.
\item LSEG, FX Venues Colo Connectivity Technical Deployment Guide, version 2.6 (June 2026).
\item Global Foreign Exchange Committee, Report on Last Look (August 2021), and the FX Global Code, Principle 17.
\end{omsources}

\section{Exercises}

\begin{exercise}[$\star$]\label{exo:nw:fx-connectivity:1}
Where does EBS match New York FX spot, and where does it match London FX spot and NDFs?
\end{exercise}

\begin{exercise}[$\star$]\label{exo:nw:fx-connectivity:2}
What does a credit-screened book show a firm, and why may its quantity mislead?
\end{exercise}

\begin{exercise}[$\star$]\label{exo:nw:fx-connectivity:3}
What is the one-way fibre floor from LD4 to TY3, and the round trip at a \omterm{def:nw:the-north-american-map:factor}{route factor} of 1.5?
\end{exercise}

\begin{exercise}[$\star\star$]\label{exo:nw:fx-connectivity:4}
Using \cref{prop:nw:fx-connectivity:hold}, why does a move in London leave the London provider's quotes only
\qty{0.1}{\milli\second} stale at every site?
\end{exercise}

\begin{exercise}[$\star\star$]\label{exo:nw:fx-connectivity:5}
A Tokyo trader sends an order for a New York-matched pair through the Tokyo gateway. How long does it take to reach the engine in the
model, and what would a New York server gain?
\end{exercise}

\begin{exercise}[$\star\star$]\label{exo:nw:fx-connectivity:6}
Why does an aggregator tend to show the stalest quote after a move?
\end{exercise}

\begin{exercise}[$\star\star\star$]\label{exo:nw:fx-connectivity:7}
\emph{Coding.} With \texttt{firm\_fxlinks.caught}, what hold catches 99\% of informed requests when the provider's information path is
at 2.46 and the model's jitter is \qty{0.5}{\milli\second}?
\end{exercise}

\begin{exercise}[$\star\star\star$]\label{exo:nw:fx-connectivity:8}
\emph{Find the flaw.} ``Our Tokyo fills are toxic because our quotes are 140 milliseconds stale there; a 150-millisecond last look
window will fix it.''
\end{exercise}

\section{Problem: USDJPY at Three Sites}

\begin{problem}[{Weekend problem --- a London provider's quotes at the Tokyo site}]\label{pb:nw:fx-connectivity:1}
A liquidity provider in London LD4 streams USDJPY to venues in London, Secaucus and Tokyo. Suppose the news that moves USDJPY
breaks in Tokyo. Use the chapter's model: every path at 1.5 times the fibre floor unless stated, \qty{0.1}{\milli\second} of pricing, and
jitter of \qty{0.5}{\milli\second} a leg where asked.

\medskip\noindent\textbf{Part I --- Distances.}
\begin{enumerate}
\item What are the one-way fibre floors from LD4 to NY5 and to TY3?
\item What are they at a \omterm{def:nw:the-north-american-map:factor}{route factor} of 1.5?
\item How stale is the provider's quote at each site after a Tokyo move?
\item And after a New York move?
\end{enumerate}

\medskip\noindent\textbf{Part II --- The price check.}
\begin{enumerate}[resume]
\item With every path at 1.5, what hold does the price check need to catch a Tokyo taker who saw a Tokyo move?
\item What share does the check catch at once, and after two milliseconds, with the jitter?
\item If the provider learns of Tokyo moves over a backbone at 2.46, what is the staleness and the hold?
\item What hold catches 99\% then?
\end{enumerate}

\medskip\noindent\textbf{Part III --- The rules.}
\begin{enumerate}[resume]
\item What does Principle 17 say the checks are for?
\item Can the provider lengthen its window to \qty{45}{\milli\second}?
\item What are its allowed remedies?
\item Where would a Tokyo pricing engine sit, and what would it change for the Tokyo site?
\end{enumerate}

\medskip\noindent\textbf{Part IV --- The verdict.}
\begin{enumerate}[resume]
\item State the \emph{named result}: the quote staleness at the Tokyo site of a London-hosted provider, and the hold time that would
  neutralise it.
\item Which numbers are published and which assumed?
\item What does an aggregator in Tokyo see after a Tokyo move?
\item How should the provider measure its staleness in production?
\item What does the Tokyo gateway mean for a Tokyo taker on EBS?
\item How would microwave paths for takers change the result?
\item Where does this go in chapter~29's plan?
\item In one sentence: what decides whether last look can protect a distant provider?
\end{enumerate}
\end{problem}

\section{Interview questions}

\begin{interviewq}[$\star$ \iqroles{developer}]\label{iq:nw:fx-connectivity:1}
Where would you put servers to trade spot FX on the main venues, and why?
\end{interviewq}

\begin{interviewq}[$\star\star$ \iqroles{trader}]\label{iq:nw:fx-connectivity:2}
What is last look for, and what is it not for?
\end{interviewq}

\begin{interviewq}[$\star\star$ \iqroles{developer, researcher}]\label{iq:nw:fx-connectivity:3}
How would you measure how stale your quotes are at each venue?
\end{interviewq}

\begin{interviewq}[$\star\star$ \iqroles{developer}]\label{iq:nw:fx-connectivity:4}
Your fills on one venue have become toxic since a network change. What do you look at?
\end{interviewq}

\begin{interviewq}[$\star\star$ \iqroles{researcher}]\label{iq:nw:fx-connectivity:5}
Why can an aggregator's best price be worse than it looks?
\end{interviewq}

\begin{interviewq}[$\star\star\star$ \iqroles{developer, researcher}]\label{iq:nw:fx-connectivity:6}
Design the pricing and streaming infrastructure of a liquidity provider quoting G10 pairs to clients in London, New York and Tokyo.
\end{interviewq}
